\documentclass[11pt,a4paper]{article}
\usepackage[T1]{fontenc}
\usepackage[utf8]{inputenc}
\IfFileExists{lmodern.sty}{\usepackage{lmodern}}{\usepackage{mathptmx}}
\usepackage[french]{babel}
\usepackage[margin=2.3cm]{geometry}
\usepackage{amsmath,amssymb,mathtools,bm}
\usepackage{graphicx,booktabs,array,multirow}
\usepackage{xcolor}
\usepackage{tikz}
\usetikzlibrary{arrows.meta,patterns,decorations.pathreplacing,calc}
\usepackage{listings}
\usepackage{tcolorbox}
\usepackage{enumitem}
\usepackage[hidelinks]{hyperref}
\usepackage{caption}
\captionsetup{font=small,labelfont=bf}
\usepackage{siunitx}
\sisetup{output-decimal-marker={,},group-separator={\,}}

\definecolor{bleu}{HTML}{2a78d6}
\definecolor{encre}{HTML}{222222}
\definecolor{fond}{HTML}{f4f7fb}
\lstset{basicstyle=\ttfamily\small,backgroundcolor=\color{fond},frame=single,rulecolor=\color{bleu!40},
  columns=fullflexible,breaklines=true,keepspaces=true,showstringspaces=false,
  literate={é}{{\'e}}1 {è}{{\`e}}1 {à}{{\`a}}1 {ê}{{\^e}}1 {ç}{{\c{c}}}1 {ô}{{\^o}}1 {É}{{\'E}}1 {â}{{\^a}}1 {î}{{\^i}}1 {ù}{{\`u}}1 {µ}{{$\mu$}}1 {«}{{\guillemotleft}}1 {»}{{\guillemotright}}1}
\newtcolorbox{encadre}[1][]{colback=fond,colframe=bleu!60,boxrule=0.5pt,arc=2pt,left=6pt,right=6pt,top=4pt,bottom=4pt,#1}

\newcommand{\dd}{\mathrm{d}}
\newcommand{\ii}{\mathrm{i}}
\newcommand{\hatt}[1]{\widehat{#1}}
\newcommand{\cs}{\texttt{chausspec}}
\renewcommand{\epsilon}{\varepsilon}

\title{\textbf{\cs}\\[4pt]\large Calcul semi-analytique spectral des chaussées multicouches\\
viscoélastiques sous empreintes hétérogènes et charges roulantes\\[8pt]
\normalsize Notice théorique, validation et mode d'emploi --- code v0.4}
\author{Lucas \textsc{David}\\ \small LTDS -- ENTPE, thèse en partenariat avec le STAC}
\date{Septembre 2026}

\begin{document}
\sloppy
\maketitle

\begin{abstract}
\cs{} calcule la réponse mécanique (déplacements, déformations, contraintes) d'une chaussée
multicouche horizontalement infinie, à couches élastiques ou viscoélastiques linéaires, soumise à
un chargement de surface de forme quelconque : empreintes uniformes, empreintes hétérogènes
analytiques, cartes de pression mesurées, atterrisseurs multi-roues, efforts tangentiels. Trois
régimes sont traités : statique, harmonique (HWD en fréquentiel) et charge roulante à vitesse
constante en régime permanent. La méthode est semi-analytique : double transformée de Fourier dans
le plan, solution exacte de l'équation de Navier dans chaque couche, assemblage par matrices de
rigidité, retour numérique dans l'espace physique. La loi 2S2P1D est utilisée telle quelle, sans
passer par un modèle rhéologique discret. Le code est validé contre des solutions analytiques
(Love), un code multicouche indépendant (PyMastic), une intégrale d'hérédité et le modèle
éléments finis COMSOL du TFE : un calcul COMSOL de \num{7.4} millions de degrés de liberté est
reproduit en quelques dizaines de secondes sur un ordinateur portable.
\end{abstract}

\tableofcontents
\newpage

%=====================================================================================================
\section{Objet et positionnement}
%=====================================================================================================
\subsection{Pourquoi ce programme}
Le TFE a montré, par un modèle éléments finis 3D (COMSOL), que la distribution de la pression de
contact sous un pneumatique d'avion modifie peu la réponse en base de couche liée, mais fortement
la réponse près de la surface. Le jury a posé deux questions de généralisation :
\begin{enumerate}[nosep]
  \item mener la même étude \emph{en domaine fréquentiel}, avec le 2S2P1D directement ;
  \item construire une solution \emph{(semi-)analytique} par superposition de solutions élémentaires.
\end{enumerate}
\cs{} répond aux deux. Une empreinte quelconque est décomposée en ondes planes $e^{\ii(k_1x+k_2y)}$
(transformée de Fourier) ; la réponse à chaque onde est calculée exactement ; la réponse totale est
la superposition. Pour une charge roulante, chaque onde plane est vue par le matériau comme une
sollicitation harmonique de pulsation $\omega=-k_1V$ : le 2S2P1D s'applique donc onde par onde, sans
approximation par un modèle de Kelvin-Voigt généralisé.

Ce programme sert en thèse à :
\begin{itemize}[nosep]
  \item vérifier rapidement des calculs éléments finis (quelques secondes à quelques minutes, contre des heures) ;
  \item mener des études paramétriques massives : vitesse, température, empreinte, structure ;
  \item accueillir directement les cartes de pression mesurées par le prototype du STAC
        (format « carte de pixels », \S\ref{sec:empreintes}) ;
  \item poser des problèmes inverses (rétrocalcul HWD, identification d'empreinte) grâce au coût réduit.
\end{itemize}

\subsection{Ce qui existe déjà, et ce qui est propre à ce travail}\label{sec:originalite}
La méthode (transformée de Fourier 2D, multicouche viscoélastique, charge roulante en régime
permanent) \textbf{n'est pas nouvelle}. Les travaux antérieurs sont nombreux et devront être cités :
\begin{itemize}[nosep]
  \item \textbf{Burmister} (1943--1945) : multicouche élastique, transformée de Hankel ; base
        d'Alizé-LCPC, de KENLAYER, de PyMastic.
  \item \textbf{ViscoRoute} (Duhamel et al., 2005 ; Chabot et al., 2010 ; Chupin et al., 2010) :
        multicouche viscoélastique (loi de Huet-Sayegh) sous charges roulantes, double transformée de
        Fourier dans le repère mobile, interfaces glissantes. Il a été conçu justement pour les
        déformations transversales asymétriques mesurées sous avion (planche PEP). Les charges y
        sont uniformes (rectangulaires ou elliptiques).
  \item \textbf{3D-Move} (Siddharthan et al., depuis 1998, Université du Nevada) : méthode des couches
        finies continues en Fourier, charges mobiles, contraintes de contact normales et
        tangentielles non uniformes de forme quelconque (dont les cartes VRSPTA de De Beer).
  \item Méthodes éléments finis 2,5D et couches finies de Fourier récentes, appliquées aux
        empreintes non uniformes de poids lourds.
\end{itemize}
À noter : le jeu de paramètres 2S2P1D du TFE a $\beta=\infty$ (pas d'amortisseur linéaire). Le
2S2P1D se réduit alors à la loi de Huet-Sayegh, celle de ViscoRoute.

\begin{encadre}
\textbf{Ce qu'on peut revendiquer honnêtement.} Il s'agit d'une implémentation propre et
validée d'une méthode connue, orientée vers un usage précis. Les apports sont :
(i) une chaîne ouverte et maîtrisée, qui accepte toute loi en module complexe (2S2P1D, KVG,
Maxwell, tabulée) ; (ii) des spectres exacts pour les empreintes hétérogènes paramétriques et pour
les cartes de pression en pixels, prévus pour les mesures du prototype STAC ; (iii) un traitement
numérique des grandes longueurs d'onde par partition de l'unité (\S\ref{sec:partition}). Il supprime
la périodisation et rend compte de la très longue mémoire des enrobés derrière une charge lente.
À ma connaissance, cette partie n'est pas documentée sous cette forme ; c'est à vérifier dans la
bibliographie. (iv) L'application aux gros porteurs et à la fissuration par le haut. L'originalité
de la thèse se situera dans (ii) et (iv), et dans ce que l'outil permettra de faire (études
paramétriques, problèmes inverses), pas dans la méthode elle-même.
\end{encadre}

\subsection{Hypothèses}
\begin{itemize}[nosep]
  \item Couches horizontales, infinies en plan, homogènes et isotropes ; petites perturbations.
  \item Comportement linéaire : élastique, ou viscoélastique linéaire décrit par un module complexe
        $E^*(\omega)$ et un coefficient de Poisson $\nu^*(\omega)$, éventuellement complexe.
  \item Quasi-statique : \textbf{inertie négligée}, comme dans le modèle COMSOL du TFE. C'est
        légitime tant que $V$ reste petit devant la célérité des ondes de cisaillement du sol
        (une centaine de m/s). C'est acceptable au roulage lent et discutable au décollage (\S\ref{sec:limites}).
  \item Température uniforme par couche. Un gradient thermique se traite en découpant l'enrobé en
        sous-couches à températures différentes.
  \item Interfaces collées ou parfaitement glissantes ; fond semi-infini ou substratum rigide
        (collé ou glissant).
  \item Charge roulante : vitesse constante, régime permanent. La charge est supposée rouler depuis
        $t=-\infty$ ; c'est l'hypothèse de ViscoRoute.
\end{itemize}

%=====================================================================================================
\section{Problème mécanique}
%=====================================================================================================
\begin{figure}[h]
\centering
\begin{tikzpicture}[scale=1.0,>=Stealth]
  \fill[blue!6] (0,0) rectangle (10,-0.8); \fill[orange!8] (0,-0.8) rectangle (10,-1.9);
  \fill[green!6] (0,-1.9) rectangle (10,-3.0);
  \fill[pattern=north east lines,pattern color=gray!60] (0,-3.0) rectangle (10,-3.3);
  \draw (0,0)--(10,0); \draw (0,-0.8)--(10,-0.8); \draw (0,-1.9)--(10,-1.9); \draw[thick] (0,-3.0)--(10,-3.0);
  \node[anchor=west] at (0.2,-0.4) {\small couche 1 : $h_1$, $E_1^*(\omega)$, $\nu_1$};
  \node[anchor=west] at (0.2,-1.35) {\small couche 2 : $h_2$, $E_2$, $\nu_2$};
  \node[anchor=west] at (0.2,-2.45) {\small couche $n$ (semi-infinie, ou reposant sur un substratum rigide)};
  \node[anchor=west] at (10.05,-0.8) {\scriptsize interface collée ou glissante};
  % charges
  \foreach \x in {3.2,3.45,...,4.4} \draw[->,thick,bleu] (\x,{0.75+0.25*sin(900*(\x-3.2))})--(\x,0.02);
  \foreach \x in {6.2,6.45,...,7.4} \draw[->,thick,bleu] (\x,{0.75+0.25*sin(900*(\x-6.2))})--(\x,0.02);
  \node[bleu] at (5.3,1.1) {\small $p(x,y)$, $q_x$, $q_y$};
  \draw[->,thick] (7.9,0.6)--(9.2,0.6) node[right] {\small $V$};
  % axes
  \draw[->] (-1.9,0)--(-0.8,0) node[right] {\small $x$};
  \draw[->] (-1.9,0)--(-1.9,-1.1) node[below] {\small $z$};
  \draw[->] (-1.9,0)--(-1.35,0.45) node[right] {\small $y$};
  \node[anchor=west] at (7.9,0.25) {\scriptsize sens de roulement};
\end{tikzpicture}
\caption{Géométrie et conventions : $x$ longitudinal (sens de roulement), $y$ transversal, $z$
vers le bas depuis la surface. \textbf{Attention} : le TFE utilise $x$ transversal et $y$
longitudinal ; $\epsilon_{xx}^{\text{TFE}}=\epsilon_{yy}$ et $\epsilon_{xz}^{\text{TFE}}=\epsilon_{yz}$ ici.}
\label{fig:geom}
\end{figure}

Dans chaque couche, l'équilibre quasi-statique et la loi de Hooke (en fréquentiel, voir \S\ref{sec:corresp}) s'écrivent
\begin{equation}
\operatorname{div}\bm\sigma=\bm 0,\qquad
\bm\sigma=\lambda\,\operatorname{tr}(\bm\epsilon)\,\bm 1+2\mu\,\bm\epsilon,\qquad
\bm\epsilon=\tfrac12(\nabla\bm u+\nabla\bm u^{\mathsf T}),
\end{equation}
soit l'équation de Navier $(\lambda+\mu)\nabla(\operatorname{div}\bm u)+\mu\Delta\bm u=\bm 0$.
Les coefficients de Lamé se déduisent de $(E,\nu)$ par
$\mu=E/[2(1+\nu)]$ et $\lambda=E\nu/[(1+\nu)(1-2\nu)]$.
La surface $z=0$ porte une pression $p(x,y)$ (positive vers le bas) et des efforts tangentiels
$q_x,q_y$ exercés par le pneu sur la chaussée. La normale sortante vaut $-\bm e_z$, d'où
\begin{equation}
\sigma_{zz}(x,y,0)=-p,\qquad \sigma_{xz}(x,y,0)=-q_x,\qquad \sigma_{yz}(x,y,0)=-q_y .
\label{eq:cl_surface}
\end{equation}
Les déformations sont tensorielles ($\epsilon_{xz}=\gamma_{xz}/2$), comptées positives en extension.

%=====================================================================================================
\section{Formulation spectrale}
%=====================================================================================================
\subsection{Transformée de Fourier}
\begin{equation}
\hatt f(k_1,k_2)=\iint_{\mathbb R^2} f(x,y)\,e^{-\ii(k_1x+k_2y)}\,\dd x\,\dd y,\qquad
f(x,y)=\frac{1}{4\pi^2}\iint_{\mathbb R^2}\hatt f(k_1,k_2)\,e^{\ii(k_1x+k_2y)}\,\dd k_1\,\dd k_2 .
\end{equation}
On note $\xi=\sqrt{k_1^2+k_2^2}$ le nombre d'onde radial et $c_1=k_1/\xi$, $c_2=k_2/\xi$ les
cosinus directeurs. La dérivation $\partial_x$ devient une multiplication par $\ii k_1$, et
$\partial_y$ par $\ii k_2$.

\subsection{Viscoélasticité : principe de correspondance et régimes}\label{sec:corresp}
Pour une sollicitation harmonique $\epsilon(t)=\operatorname{Re}[\epsilon_0e^{\ii\omega t}]$, un
matériau viscoélastique linéaire répond par $\sigma(t)=\operatorname{Re}[E^*(\omega)\epsilon_0e^{\ii\omega t}]$.
Le signal temporel étant réel, $E^*(-\omega)=\overline{E^*(\omega)}$.
Si chaque composante du problème évolue en $e^{\ii\omega t}$, l'équation de Navier est inchangée,
avec $\lambda^*(\omega)$ et $\mu^*(\omega)$ complexes : c'est le principe de correspondance. Les
trois régimes de \cs{} ne diffèrent que par la pulsation associée à chaque onde plane :
\begin{center}
\begin{tabular}{lll}
\toprule
Régime & Chargement & Pulsation vue par le matériau\\
\midrule
\texttt{Static()} & $p(x,y)$ appliqué depuis toujours & $\omega=0$ (module statique $E_{00}$)\\
\texttt{Harmonic(f)} & $p(x,y)\,e^{\ii 2\pi f t}$ (HWD) & $\omega=2\pi f$ pour tout $\bm k$\\
\texttt{Moving(V)} & $p(x-Vt,y)$, régime permanent & $\omega=-k_1V$\\
\bottomrule
\end{tabular}
\end{center}
\textbf{Charge roulante.} En régime permanent, les champs sont stationnaires dans le repère
mobile $X=x-Vt$ : $\bm u(x,y,z,t)=\bm U(x-Vt,y,z)$. La composante
$\hatt{\bm U}(k_1,k_2,z)\,e^{\ii(k_1X+k_2y)}$ vaut, en un point matériel fixe,
$\hatt{\bm U}\,e^{\ii(k_1x+k_2y)}e^{-\ii k_1Vt}$ : c'est une sollicitation harmonique de pulsation
$\omega=-k_1V$. Chaque onde plane est donc un problème élastique à modules complexes
$\lambda^*(-k_1V)$, $\mu^*(-k_1V)$. Le signe fixe le sens du retard : la réponse est décalée
\emph{derrière} la charge ($X<0$), ce que vérifie la validation V4 (figure~\ref{fig:v4}).
Les ondes courtes ($|k_1|$ grand) voient le module vitreux ; les ondes longues voient le module
statique. Une empreinte de \SI{0.5}{m} à \SI{0.66}{m/s} mobilise donc une large bande de
fréquences, de $10^{-4}$ à quelques dizaines de hertz, et pas « la » fréquence de \SI{1}{Hz} de la
règle de Perret. C'est exactement la limite du module équivalent discutée au \S3.6 du TFE.

\subsection{Lois de comportement disponibles}
\paragraph{2S2P1D} (Olard et Di Benedetto, 2003) :
\begin{equation}
E^*(\omega)=E_{00}+\frac{E_0-E_{00}}{1+\delta(\ii\omega\tau)^{-k}+(\ii\omega\tau)^{-h}+(\ii\omega\beta\tau)^{-1}},
\qquad \tau(T)=a_T(T)\,\tau(T_{\text{ref}}),
\end{equation}
avec, en option, le décalage WLF $\log_{10}a_T=-C_1(T-T_{\text{ref}})/(C_2+T-T_{\text{ref}})$ et
un coefficient de Poisson complexe
$\nu^*=\nu_{00}+(\nu_0-\nu_{00})(E^*-E_{00})/(E_0-E_{00})$. Les puissances complexes
$(\ii\omega\tau)^{-k}$ sont prises sur la détermination principale.

\paragraph{Kelvin-Voigt généralisé (KVG)} (forme du TFE, Tableau 7) :
$1/E^*(\omega)=1/E_0+\sum_i 1/[E_i(1+\ii\omega\tau_i)]$.
\paragraph{Maxwell généralisé (Prony)} : $E^*(\omega)=E_\infty+\sum_iE_i\,\ii\omega\tau_i/(1+\ii\omega\tau_i)$.
\paragraph{Module figé} : élastique de module $|E^*(f)|$ d'une autre loi (approche « module
équivalent »).

\begin{figure}[h]
\centering\includegraphics[width=0.92\linewidth]{figures/courbes_maitresses.pdf}
\caption{BB-GB du TFE à \SI{9.3}{\celsius} : 2S2P1D du Tableau 6 ($\tau=\SI{1.22}{s}$) et KVG
à 9 branches du Tableau 7. Le KVG approche le 2S2P1D à mieux que 10\,\% sur la plage utile. À
\SI{1}{Hz}, $|E^*|$ vaut \SI{14877}{MPa} (2S2P1D) et \SI{14482}{MPa} (KVG), et non
\SI{11670}{MPa}. Cette dernière valeur correspond à environ \SI{0.3}{Hz}, ou à $\tau\approx\SI{0.35}{s}$
à \SI{1}{Hz}. Le point est à trancher avant l'article MAIREINFRA.}
\label{fig:master}
\end{figure}

\subsection{Découplage P-SV / SH}
Pour des couches isotropes, on cherche les transformées sous la forme
\begin{equation}
\hatt u_x=\frac{\ii}{\xi}\big(k_1U-k_2V\big),\qquad
\hatt u_y=\frac{\ii}{\xi}\big(k_2U+k_1V\big),\qquad
\hatt u_z=W,
\label{eq:ansatz}
\end{equation}
où $U(z)$, $W(z)$ et $V(z)$ ne dépendent que de $\xi$, de $z$ et des modules. $U$ porte la partie
irrotationnelle du déplacement horizontal, $V$ sa partie solénoïdale. Avec $\operatorname{div}\hatt{\bm u}=W'-\xi U$,
l'équation de Navier se scinde en un système « P-SV » et une équation « SH » :
\begin{align}
&(\lambda+\mu)\,\xi\,(W'-\xi U)+\mu\,(U''-\xi^2U)=0, \label{eq:psv1}\\
&(\lambda+\mu)\,(W''-\xi U')+\mu\,(W''-\xi^2W)=0, \label{eq:psv2}\\
&\mu\,(V''-\xi^2V)=0 .
\end{align}
Les contraintes utiles aux interfaces s'expriment par
\begin{equation}
\hatt\sigma_{zz}=S=(\lambda+2\mu)W'-\lambda\xi U,\qquad
\hatt\sigma_{xz}=\frac{\ii}{\xi}(k_1T-k_2R),\qquad
\hatt\sigma_{yz}=\frac{\ii}{\xi}(k_2T+k_1R),
\end{equation}
avec $T=\mu(U'+\xi W)$ et $R=\mu V'$. Sous chargement purement normal, $V\equiv0$ : seul le
problème P-SV est sollicité.

\subsection{Solution générale dans une couche}
On pose $\kappa=(\lambda+3\mu)/(\lambda+\mu)=3-4\nu$ (complexe en viscoélasticité). Dans une couche
d'épaisseur $h$ et de cote locale $s\in[0,h]$, on note $\tau=\xi s$ et $t=\xi(h-s)$. La solution
générale de \eqref{eq:psv1}--\eqref{eq:psv2} (obtenue avec SymPy, annexe~\ref{ann:sympy}) s'écrit
\begin{equation}
\boxed{\begin{aligned}
W&=(A+B\tau)\,e^{-\tau}+(C+D\,t)\,e^{-t},\\
U&=(-A+\kappa B-B\tau)\,e^{-\tau}+(C-\kappa D+D\,t)\,e^{-t},\\
V&=a\,e^{-\tau}+b\,e^{-t}.
\end{aligned}}
\label{eq:solgen}
\end{equation}
La famille $(A,B)$ décroît depuis le toit de la couche, la famille $(C,D)$ depuis sa base. Toutes
les exponentielles sont bornées par 1. Contrairement à la forme classique en $e^{\pm\xi z}$ de
Burmister, il n'y a donc ni dépassement de capacité ni perte de précision pour les grands $\xi h$.
Les dérivées s'obtiennent par $\dd/\dd z=\xi\,\dd/\dd\tau=-\xi\,\dd/\dd t$ :
\begin{align*}
\tfrac{\dd W}{\dd\tau}&=(B-A-B\tau)e^{-\tau}+(C-D+Dt)e^{-t}, &
\tfrac{\dd U}{\dd\tau}&=(A-(1+\kappa)B+B\tau)e^{-\tau}+(C-(1+\kappa)D+Dt)e^{-t}.
\end{align*}
Pour la dernière couche d'un massif semi-infini, $C=D=b=0$.

\paragraph{Mise à l'échelle.} Les inconnues calculées sont les coefficients physiques multipliés
par $\mu_{\text{ref}}\,\xi$ (avec $\mu_{\text{ref}}=|\mu_1|$). Les contraintes $T,S$ sont alors en
Pa et les déplacements valent $(\cdots)/(\mu_{\text{ref}}\xi)$. Toutes les matrices sont sans
dimension et d'ordre 1.

\subsection{Conditions aux limites et raccordements}
\paragraph{Surface.} D'après \eqref{eq:cl_surface} et \eqref{eq:ansatz} :
\begin{equation}
S(0)=-\hatt p,\qquad T(0)=\ii\,(c_1\hatt q_x+c_2\hatt q_y),\qquad R(0)=\ii\,(-c_2\hatt q_x+c_1\hatt q_y).
\end{equation}
Le noyau est résolu pour trois chargements unitaires (pression normale, P-SV tangentiel, SH).
La réponse à un chargement quelconque est la combinaison linéaire de ces trois cas, pondérés par
$\hatt p$, $T(0)$ et $R(0)$.
\paragraph{Interfaces.} \emph{Collée} : continuité de $U,W,V,T,S,R$.
\emph{Glissante} (glissement parfait) : continuité de $W$ et $S$, et $T=R=0$ de part et d'autre.
\paragraph{Fond.} \texttt{halfspace} (massif semi-infini), \texttt{rigid\_bonded} ($U=W=V=0$),
\texttt{rigid\_smooth} ($W=0$, $T=R=0$). La dernière condition est celle du modèle COMSOL du TFE
(déplacement vertical bloqué sous le substratum).

\subsection{Résolution par matrices de rigidité de couche}
Pour une couche finie $j$, notons $\bm d=(U_t,W_t,U_b,W_b)$ les déplacements de ses faces
supérieure et inférieure, et $\bm s=(T_t,S_t,T_b,S_b)$ les contraintes correspondantes. En
évaluant \eqref{eq:solgen} en $s=0$ et $s=h$, on obtient $\bm d=\bm D_j\bm c$ et $\bm s=\bm Q_j\bm c$,
avec $\bm c=(A,B,C,D)$, d'où la \emph{matrice de rigidité de couche}
\begin{equation}
\bm K_j=\bm Q_j\bm D_j^{-1}\quad(4\times4),\qquad
\bm D_j^{-1}\ \text{par complément de Schur sur les blocs } 2\times2,
\end{equation}
ce qui est licite car les blocs diagonaux ont pour déterminants $-\kappa$ et $\kappa$. Les
tractions de face sont $\operatorname{diag}(-1,-1,1,1)\,\bm K_j\bm d$. L'équilibre de chaque face
(somme des tractions des deux couches adjacentes nulle, traction imposée en surface) conduit à un
système \textbf{tridiagonal par blocs $2\times2$} sur les déplacements des faces. On le résout par
l'algorithme de Thomas, vectorisé sur tous les nombres d'onde. Une interface glissante est traitée
par \emph{condensation statique} : le déplacement horizontal de la face est éliminé grâce à $T=0$,
puis remplacé par une inconnue muette. Le problème SH conduit à un système tridiagonal scalaire.
Les coefficients $\bm c$ de la couche qui contient la cote demandée se déduisent enfin des
déplacements de ses faces.

Une résolution « dense » (un seul système de $4n$ inconnues) est conservée comme référence. Les
deux méthodes coïncident à $10^{-7}$ près sur toutes les combinaisons de fonds et d'interfaces
testées (\texttt{tests/test\_kernels.py}). La méthode par rigidités est 5 à 10 fois plus rapide.

\subsection{Reconstruction des champs}
Avec $U,W,V$ et leurs dérivées à la cote $z$ :
\begin{align}
\hatt u_x&=\ii\,(c_1U-c_2V), & \hatt u_y&=\ii\,(c_2U+c_1V), & \hatt u_z&=W,\\
\hatt\epsilon_{xx}&=-\xi\,c_1(c_1U-c_2V), & \hatt\epsilon_{yy}&=-\xi\,c_2(c_2U+c_1V), & \hatt\epsilon_{zz}&=W',\\
\hatt\epsilon_{xy}&=-\xi\big[c_1c_2U+\tfrac12(c_1^2-c_2^2)V\big], &
\hatt\epsilon_{xz}&=\tfrac{\ii}{2}\big[c_1(U'+\xi W)-c_2V'\big], &
\hatt\epsilon_{yz}&=\tfrac{\ii}{2}\big[c_2(U'+\xi W)+c_1V'\big],
\end{align}
puis $\hatt{\bm\sigma}=\lambda\operatorname{tr}\hatt{\bm\epsilon}\,\bm1+2\mu\hatt{\bm\epsilon}$,
le produit étant exact dans le domaine spectral pour la viscoélasticité. Quand $\xi\to0$, les
déformations et contraintes ont des limites finies, qui dépendent de la direction. Le déplacement
d'un massif semi-infini diverge en revanche comme $1/\xi$ : la singularité est intégrable en 2D,
mais pas sur une seule ligne $k_2=0$. Le \S\ref{sec:partition} en tient compte.

\begin{encadre}
\textbf{Conséquence directe : $\epsilon_{xz}=\epsilon_{yz}=0$ en surface} dès que les efforts
tangentiels sont nuls, puisque $\sigma_{xz}(z=0)=-q_x=0$ et que $\epsilon_{xz}=\sigma_{xz}/2\mu$.
Les « cisaillements en surface » non nuls d'un calcul éléments finis sous chargement purement
vertical sont donc des artefacts d'extrapolation nodale. La grandeur pertinente est le maximum
atteint \emph{sous} la surface. Le TFE le situe vers 50 à \SI{160}{mm} ; voir \S\ref{sec:v5}.
\end{encadre}

%=====================================================================================================
\section{Retour dans l'espace physique}
%=====================================================================================================
\subsection{Spectres des empreintes}\label{sec:empreintes}
Chaque roue est une empreinte centrée en $(x_0,y_0)$ ; son spectre est multiplié par
$e^{-\ii(k_1x_0+k_2y_0)}$ et les roues se superposent. Les spectres sont \emph{exacts} (aucun
échantillonnage de la charge) :
\begin{center}\small
\begin{tabular}{p{4.3cm}p{9.6cm}}
\toprule
Empreinte & Spectre $\hatt p(k_1,k_2)$\\
\midrule
Rectangle uniforme $\ell_x\times\ell_y$ & $p\,\ell_x\operatorname{sinc}(k_1\ell_x/2)\,\ell_y\operatorname{sinc}(k_2\ell_y/2)$\\
Disque uniforme de rayon $R$ & $p\,2\pi R^2\,J_1(\xi R)/(\xi R)$\\
Séparable $A\,f_x(x)\,f_y(y)$ & $A\,\hatt f_x(k_1)\,\hatt f_y(k_2)$\\
\quad demi-ellipse $\sqrt{1-(s/c)^2}$ & $\pi c\,J_1(kc)/(kc)$\\
\quad paires de gaussiennes (Annexe III du TFE) & $\sum_iP_i\,\sigma_i\sqrt{2\pi}\,e^{-k^2\sigma_i^2/2}\,2\cos(k s_i)$\\
\quad profil tabulé, interpolé linéairement & $d\operatorname{sinc}^2(kd/2)\sum_nf_n e^{-\ii ks_n}$\\
Carte de pixels $P_{ji}$ (pas $d_x,d_y$) & $d_x\operatorname{sinc}(k_1d_x/2)\,d_y\operatorname{sinc}(k_2d_y/2)\sum_{j,i}P_{ji}e^{-\ii(k_1x_i+k_2y_j)}$\\
\bottomrule
\end{tabular}
\end{center}
La carte de pixels est l'entrée prévue pour les mesures du prototype STAC : la transformée de la
fonction constante par pixel est calculée exactement, par produits matriciels séparables.
La résultante vaut toujours $\hatt p(0,0)$ ; la méthode \texttt{scaled\_to(F)} normalise une
empreinte à une charge donnée (le $K_{\text{charge}}$ du TFE).

\begin{figure}[h]
\centering\includegraphics[width=\linewidth]{figures/empreintes.pdf}
\caption{Les trois empreintes du TFE à \SI{370}{kN}. L'intégrale exacte du profil de l'Annexe III
vaut \SI{142.1}{kN} avant mise à l'échelle (et non 141,2), d'où $K_{\text{charge}}=2{,}603$ et
$p_{\max}=\SI{3.44}{MPa}$.}
\label{fig:empreintes}
\end{figure}

\subsection{Discrétisation}
Le champ est calculé sur une grille $N_x\times N_y$ de pas $d_x=L_x/N_x$, $d_y=L_y/N_y$ ; les
nombres d'onde sont ceux de la FFT ($\Delta k_1=2\pi/L_x$). Deux règles :
\begin{itemize}[nosep]
\item \textbf{résolution} : le pas doit résoudre les plus fins détails de la charge (nervures :
      $\sigma_i\approx\SI{17}{mm}$, d'où $d\lesssim\SI{8}{mm}$) pour les grandeurs proches de la
      surface. En profondeur, le facteur $e^{-\xi z}$ filtre naturellement les ondes courtes ;
\item \textbf{domaine} : grâce au \S\ref{sec:partition}, quelques fois l'emprise du chargement
      suffisent. Pour le bogie d'A340, $\SI{16}{m}\times\SI{8}{m}$ convient.
\end{itemize}
Un filtre gaussien optionnel (\texttt{filter\_width}) lisse la charge à l'échelle du maillage,
comme le fait implicitement un modèle éléments finis. Sous une charge en créneau, les valeurs
exactement en surface convergent lentement (phénomène de Gibbs) : elles sont d'ailleurs
singulières en théorie au bord d'un créneau.

\subsection{Traitement des grandes longueurs d'onde : partition de l'unité}\label{sec:partition}
\paragraph{Le problème.} Une FFT seule calcule la réponse à un réseau \emph{périodique} de
chargements. C'est gênant dans trois cas :
(a) en charge roulante viscoélastique, un point matériel garde la mémoire de la charge pendant
$\tau_{\max}$. Le KVG du TFE a des branches jusqu'à \SI{1.6e4}{s}, soit une traînée de
$V\tau_{\max}\approx\SI{10}{km}$ derrière la charge. Le noyau varie sur une échelle
\emph{logarithmique} en $k_1$ (figure~\ref{fig:partition}), très en dessous du pas $\Delta k_1$ ;
(b) les champs décroissent lentement en profondeur (en $r^{-3}$ pour $\sigma_{zz}$ sur un sol
mou), et les images périodiques se cumulent ;
(c) la déflexion d'un massif semi-infini est singulière en $1/\xi$.
Le premier défaut a été mis en évidence en validation (V4) : écarts de 5 à 20\,\% avec un
traitement naïf de la colonne $k_1=0$.

\paragraph{La solution.} On scinde le spectre avec deux gaussiennes
$\varphi(k_1)=e^{-k_1^2/2(b\Delta k_1)^2}$ et $\psi(k_2)=e^{-k_2^2/2(b\Delta k_2)^2}$ ($b=2$ par
défaut) :
\begin{equation}
F=\underbrace{(1-\varphi)F}_{\text{FFT 2D}}
+\underbrace{\varphi(1-\psi)F}_{\substack{\text{intégrale continue en }k_1\\\text{FFT en }y}}
+\underbrace{\varphi\psi F}_{\text{intégrale continue 2D}} .
\end{equation}
\begin{itemize}[nosep]
\item La partie FFT est de moyenne nulle en $x$ ($1-\varphi(0)=0$) et régulière : son image dans
      l'espace physique est localisée, et la périodisation y est inoffensive.
\item Les parties basses sont intégrées \emph{continûment}, par une quadrature de Gauss-Legendre sur
      des panneaux géométriques vers $k=0$ (rapport 4, jusqu'à $10^{-7}$ fois la largeur de bande)
      puis réguliers. Les phases $e^{\ii k_1X}$ sont calculées exactement aux points de sortie.
      Ces parties ne sont \emph{pas} périodiques.
\item La singularité en $1/\xi$ du déplacement est intégrable en 2D : la déflexion absolue d'un
      massif semi-infini devient calculable (validée au \S\ref{sec:v1}).
\end{itemize}
Effet mesuré sur le cas du TFE ($\epsilon_T$ en base, KVG, empreinte rectangulaire) :
\begin{center}\small
\begin{tabular}{lcccccc}
\toprule
Domaine $L_x\times L_y$ (m) & $16\times8$ & $32\times8$ & $64\times8$ & $32\times16$ & $32\times32$ & $64\times64$\\
\midrule
FFT seule + moyenne de cellule (µdef) & 341,1 & 339,7 & 338,0 & 337,1 & 336,2 & 336,0\\
Partition de l'unité (µdef) & \textbf{335,9} & --- & --- & 336,0 & 336,0 & 336,0\\
\bottomrule
\end{tabular}
\end{center}

\begin{figure}[h]
\centering\includegraphics[width=0.78\linewidth]{figures/partition.pdf}
\caption{Module vu par une onde plane $e^{\ii k_1X}$ sous charge roulante (KVG, $V=\SI{0.66}{m/s}$).
Le module tombe à $E_{00}=\SI{65}{MPa}$ pour $k_1\lesssim10^{-4}$, très en dessous du premier point de
FFT : la zone bleutée est intégrée continûment.}
\label{fig:partition}
\end{figure}

\subsection{Solveur axisymétrique (Hankel)}
Pour une charge circulaire uniforme, en statique ou en harmonique (plaque HWD), le même noyau est
inversé par transformée de Hankel :
$f(r)=\frac1{2\pi}\int_0^\infty F(\xi)J_0(\xi r)\xi\,\dd\xi$,
$u_r=-\frac1{2\pi}\int_0^\infty U(\xi)J_1(\xi r)\xi\,\dd\xi$.
L'intégrale est calculée par une quadrature de Gauss sur des panneaux adaptés aux oscillations.
Ce solveur est très rapide (\SI{0.1}{s}) et sert à la validation croisée.

\subsection{Post-traitements}
Déformations principales (valeurs propres du tenseur $3\times3$ en chaque point), maxima et leur
position, coupes, interpolation, et \textbf{signal temporel d'une jauge fixe} :
en charge roulante, une jauge en $(x_g,y_g)$ voit $f(X=x_g-Vt,\,y_g)$.
Sur une interface, l'option \texttt{side="above"} (défaut) prend la couche supérieure, par exemple
pour la base de la couche liée, et \texttt{side="below"} la couche inférieure.

\subsection{Coût}
Le coût est d'environ \SI{15}{\micro s} par nombre d'onde pour une structure à 5 couches (Python
et NumPy, un c\oe ur). Un calcul du bogie d'A340 à \SI{8}{mm} de résolution
($2048\times2048$) prend environ une minute sur un portable, contre plusieurs heures pour le
modèle COMSOL de \num{7.4} millions de ddl. Pour un chargement symétrique en $y$,
\texttt{sym\_y=True} divise le temps par deux.

%=====================================================================================================
\section{Validation}
%=====================================================================================================
Les scripts sont dans \texttt{validation/}. Chacun produit un fichier CSV de résultats.

\subsection{V1 --- Massif semi-infini élastique, charge circulaire (Love)}\label{sec:v1}
$E=\SI{100}{MPa}$, $\nu=0{,}3$, $p=\SI{0.7}{MPa}$, $a=\SI{0.15}{m}$. Sur l'axe, les solutions exactes
sont $\sigma_{zz}=-p[1-z^3/R^3]$,
$\sigma_{rr}=-\frac p2[(1+2\nu)-2(1+\nu)z/R+z^3/R^3]$ et
$w=\frac{(1+\nu)pa}{E}[a/R+(1-2\nu)(R-z)/a]$, avec $R=\sqrt{a^2+z^2}$.
\begin{figure}[h]
\centering\includegraphics[width=0.9\linewidth]{figures/v1_love.pdf}
\caption{V1 : accord à $10^{-4}$ près en profondeur, pour les contraintes comme pour la déflexion
absolue (FFT 2D incluse, grâce au \S\ref{sec:partition}). En $z=0$, sous le créneau, les deux
méthodes numériques tronquent le spectre : c'est le phénomène de Gibbs.}
\end{figure}

\subsection{V2 --- Multicouche élastique : comparaison à PyMastic}
Structure à 3 couches ($E$ = 5000 / 200 / \SI{50}{MPa}, $h$ = 0,20 / \SI{0.30}{m} / semi-infini,
$\nu$ = 0,35 / 0,35 / 0,40), charge circulaire $p=\SI{0.7}{MPa}$, $a=\SI{0.15}{m}$, interfaces collées.
PyMastic (Nakhaei, 2020) est un code Burmister indépendant, validé contre KENPAVE.
\begin{center}\small
\begin{tabular}{rr|rr|rr|rr}
\toprule
$z$ (m) & $r$ (m) & \multicolumn{2}{c|}{$w$ (mm)} & \multicolumn{2}{c|}{$\sigma_{z}$ (MPa, compr. $>0$)} & \multicolumn{2}{c}{$\epsilon_r$ (µdef)}\\
 & & PyMastic & \cs & PyMastic & \cs & PyMastic & \cs\\
\midrule
0,1000 & 0,0 & 0,5249 & 0,5250 & 0,4229 & 0,4229 & 17,17 & 17,18\\
0,1000 & 0,3 & 0,4524 & 0,4525 & 0,0114 & 0,0114 & -10,27 & -10,27\\
0,1999 & 0,0 & 0,5140 & 0,5141 & 0,0732 & 0,0732 & 149,86 & 149,86\\
0,1999 & 0,3 & 0,4498 & 0,4499 & 0,0288 & 0,0288 & 11,61 & 11,61\\
0,2001 & 0,0 & 0,5140 & 0,5140 & 0,0731 & 0,0731 & 150,01 & 150,01\\
0,2001 & 0,3 & 0,4498 & 0,4499 & 0,0288 & 0,0288 & 11,64 & 11,65\\
0,3500 & 0,0 & 0,4693 & 0,4693 & 0,0350 & 0,0350 & 117,35 & 117,36\\
0,3500 & 0,3 & 0,4284 & 0,4285 & 0,0210 & 0,0210 & 54,22 & 54,23\\
0,4999 & 0,0 & 0,4366 & 0,4367 & 0,0194 & 0,0194 & 151,48 & 151,48\\
0,4999 & 0,3 & 0,4057 & 0,4058 & 0,0153 & 0,0153 & 88,93 & 88,93\\
0,5001 & 0,0 & 0,4365 & 0,4366 & 0,0194 & 0,0194 & 151,50 & 151,50\\
0,5001 & 0,3 & 0,4057 & 0,4058 & 0,0153 & 0,0153 & 88,95 & 88,95\\
0,8000 & 0,0 & 0,3476 & 0,3477 & 0,0118 & 0,0118 & 91,72 & 91,73\\
0,8000 & 0,3 & 0,3334 & 0,3334 & 0,0104 & 0,0104 & 71,30 & 71,31\\
\bottomrule
\end{tabular}
\end{center}
Accord à 4 chiffres significatifs partout, de part et d'autre des interfaces. En surface
($z=0$, non montré), les deux codes tronquent le spectre du créneau différemment. En interfaces
glissantes, PyMastic diverge numériquement dans les couches supérieures ; seule la couche de fond,
qui coïncide, est exploitable. D'où V3.

\subsection{V3 --- Interfaces glissantes}
Référence : la même structure, où chaque interface glissante est remplacée par une couche
intercalaire collée d'épaisseur $e$, quasi incompressible ($\nu=0{,}5-5\cdot10^{-11}$) et de
module $E=0{,}05\,e$ (MPa, $e$ en m). Sa raideur en cisaillement s'annule donc quand $e\to0$.
\begin{center}\small
\begin{tabular}{lcccc}
\toprule
$e$ (mm) & 4 & 1 & 0,25 & 0,0625\\
\midrule
écart max. sur $w$ & 3,3e-02 & 5,4e-03 & 6,9e-04 & 7,0e-05\\
écart max. sur $\sigma_r$ & 3,2e-02 & 6,9e-03 & 5,8e-04 & 1,4e-04\\
\bottomrule
\end{tabular}
\end{center}
La convergence vers la solution à glissement parfait est d'ordre 1,5 à 2 en $e$.

\subsection{V4 --- Charge roulante viscoélastique : intégrale d'hérédité}
Pour un corps homogène à coefficient de Poisson constant chargé en effort, le champ de contraintes
ne dépend pas du module. La déformation viscoélastique s'obtient donc à partir de la déformation
élastique de module unité par l'intégrale de Boltzmann
$\epsilon(t)=\int_{-\infty}^tJ(t-s)\,\dd\epsilon_{\text{él}}^{E=1}(s)$, avec la complaisance de
fluage du KVG $J(t)=1/E_0+\sum_i(1-e^{-t/\tau_i})/E_i$. Cette référence est calculée de façon
totalement indépendante du passage en fréquentiel ($\omega=-k_1V$). Cas : massif semi-infini en
KVG du TFE, disque de \SI{0.267}{m} à \SI{1.65}{MPa}, $V=\SI{0.66}{m/s}$, historique de \SI{128}{m}.
\begin{center}\small
\begin{tabular}{ll|rr|rr|r}
\toprule
$z$ (m) & comp. & max réf. (µdef) & $X$ (m) & max \cs{} (µdef) & $X$ (m) & écart max (\% du max)\\
\midrule
0,05 & $e_{xx}$ & 18,9 & 0,297 & 18,9 & 0,297 & 0,75\\
0,05 & $e_{yy}$ & 9,2 & -0,062 & 9,1 & -0,062 & 0,39\\
0,05 & $e_{zz}$ & 113,9 & -0,203 & 113,9 & -0,219 & 0,17\\
0,15 & $e_{xx}$ & 21,0 & -0,125 & 21,0 & -0,125 & 0,20\\
0,15 & $e_{yy}$ & 28,0 & -0,188 & 28,0 & -0,188 & 0,08\\
0,15 & $e_{zz}$ & 120,1 & -0,125 & 120,1 & -0,125 & 0,04\\
\bottomrule
\end{tabular}
\end{center}
Les écarts résiduels (< 0,8\,\%) viennent de la référence : son historique est tronqué et elle
converge vers \cs{} quand on l'allonge (écarts maximaux de 5\,\% avec \SI{64}{m} d'historique, 1\,\% avec \SI{128}{m},
0,2\,\% avec \SI{256}{m}). Les pics se situent derrière la charge ($X<0$) : le signe de $\omega=-k_1V$
est validé.
\begin{figure}[h]
\centering\includegraphics[width=0.8\linewidth]{figures/v4_heredite.pdf}
\caption{V4 : signal $\epsilon_{xx}(X)$ à $z=\SI{0.15}{m}$. La réponse viscoélastique est plus
grande et asymétrique ; les deux calculs se superposent.}
\label{fig:v4}
\end{figure}

\subsection{V5 --- Reproduction du modèle COMSOL du TFE}\label{sec:v5}
Structure PEP, bogie d'A340 (4 $\times$ \SI{370}{kN}, entraxes \SI{1.98}{m} $\times$ \SI{1.40}{m}),
$V=\SI{0.66}{m/s}$, \SI{9.3}{\celsius}, fond « déplacement vertical bloqué ». Le modèle COMSOL
compte \num{7.4} millions de ddl et 200 pas de temps. \cs{} calcule chaque cas en environ
\SI{2}{min} (base : $\SI{16}{m}\times\SI{8}{m}$, $512\times256$ ; surface et profils :
$2048\times2048$, pas de \SI{8}{mm} $\times$ \SI{4}{mm}). Les lois « KVG » et « élastique » sont
exactement celles du TFE ; la ligne « 2S2P1D » utilise le Tableau 6 directement.
\begin{center}\small
\begin{tabular}{ll|rr|rr}
\toprule
& & \multicolumn{2}{c|}{$\epsilon_T$ max en base (µdef)} & \multicolumn{2}{c}{$\epsilon_1$ max en surface (µdef)}\\
Loi & Empreinte & COMSOL (TFE) & \cs & COMSOL (TFE) & \cs\\
\midrule
élastique & rectangulaire & 210,0 & 211,1 (+0,5\,\%) & 155,0 & 157,7 (+1,7\,\%)\\
 & circulaire & 202,2 & 203,4 (+0,6\,\%) & 156,8 & 149,7 (-4,5\,\%)\\
 & hétérogène & 212,8 & 209,5 (-1,5\,\%) & 163,4 & 149,9 (-8,3\,\%)\\
\midrule
KVG & rectangulaire & 331,1 & 336,0 (+1,5\,\%) & 195,3 & 196,3 (+0,5\,\%)\\
 & circulaire & 318,0 & 321,5 (+1,1\,\%) & 181,5 & 172,9 (-4,7\,\%)\\
 & hétérogène & 336,6 & 329,7 (-2,1\,\%) & 202,1 & 181,1 (-10,4\,\%)\\
\midrule
2S2P1D & rectangulaire & --- & 316,1 & --- & 188,4\\
 & circulaire & --- & 302,8 & --- & 167,8\\
 & hétérogène & --- & 310,5 & --- & 174,6\\
\midrule
\multicolumn{2}{l|}{Mesure PEP (Petitjean et al.)} & \multicolumn{2}{c|}{340,1} & & \\
\bottomrule
\end{tabular}
\end{center}
\begin{figure}[h]
\centering\includegraphics[width=0.85\linewidth]{figures/v5_signaux.pdf}
\caption{V5 : déformation transversale en base de couche liée vue par une jauge sur l'axe des
roues (équivalent de la figure 3-11 du TFE). On retrouve l'asymétrie et la roue arrière plus
sollicitée en viscoélasticité, ainsi que la longue traînée après le passage.}
\end{figure}

\paragraph{Lecture.}
\begin{itemize}[nosep]
\item \textbf{En base de couche liée}, l'accord est de 0,5 à 1,5\,\% pour les empreintes uniformes,
      en élastique comme en viscoélastique. Le calcul spectral est en régime permanent (charge
      venue de $-\infty$), le calcul COMSOL part du repos 20 s plus tôt dans un domaine borné ;
      l'écart restant est de cet ordre.
\item \textbf{En surface}, l'empreinte rectangulaire, alignée sur le maillage COMSOL, est
      retrouvée à 1\,\% près. Pour le disque ($-5$\,\%) et l'empreinte hétérogène ($-8$ à
      $-10$\,\%), le calcul spectral est convergé en résolution : ses valeurs ne bougent pas quand on
      divise le pas par deux. Les écarts viennent vraisemblablement du maillage COMSOL. Des éléments
      de \SI{17}{mm} environ représentent un bord circulaire en escalier et des nervures de largeur
      $\sigma\approx\SI{17}{mm}$ avec seulement un ou deux éléments.
\item \textbf{Empreinte hétérogène en base} : $-2$\,\% (élastique) et $-2$\,\% (KVG) par rapport à
      COMSOL, alors que les empreintes uniformes sont à $+0{,}5$ à $+1{,}5$\,\%. L'empreinte simulée
      par COMSOL n'était sans doute pas exactement celle de l'Annexe III : le TFE mentionne aussi
      $p_{\max}$ = 3,51 puis \SI{3.31}{MPa}, et un facteur $370/372{,}8$, alors que l'intégrale
      exacte de l'Annexe III donne \SI{142.1}{kN}. \emph{À vérifier dans le fichier COMSOL.}
\item \textbf{2S2P1D contre KVG} : le 2S2P1D donne des déformations plus faibles d'environ 6\,\%,
      car le KVG est un peu plus souple entre \SI{0.01}{Hz} et \SI{3}{Hz} (figure~\ref{fig:master}).
      L'approximation par spectre discret n'est donc pas neutre au niveau de précision recherché.
\end{itemize}

\paragraph{Cisaillement $\epsilon_{xz}^{\text{TFE}}=\epsilon_{yz}$.} Il est rigoureusement nul en
surface (\S3.8). Son maximum est atteint en profondeur, près de l'épaulement extérieur du
pneumatique :
\begin{center}\small
\begin{tabular}{l|rr|rr}
\toprule
& \multicolumn{2}{c|}{COMSOL (TFE, Tab. 13)} & \multicolumn{2}{c}{\cs{} (KVG)}\\
Empreinte & « surface » (µdef) & $z$ du max & max (µdef) & $z$ du max\\
\midrule
rectangulaire & 143,2 & $\approx$ 160 mm & 152,4 & 120 mm\\
circulaire & 120,9 & $\approx$ 160 mm & 146,8 & 140 mm\\
hétérogène & 177,3 & $\approx$ 50 mm & 164,2 & 100 mm\\
\bottomrule
\end{tabular}
\end{center}
Le classement est conservé (hétérogène > rectangulaire > circulaire) et les ordres de grandeur sont
comparables. Mais l'écart hétérogène/rectangulaire tombe de +24\,\% à +8\,\% une fois les maxima
pris là où ils se trouvent vraiment, et non en surface. \textbf{Pour l'article, il faut retenir les
maxima en profondeur ; les valeurs « en surface » du Tableau 13 sont des artefacts.}
\begin{figure}[h]
\centering\includegraphics[width=0.8\linewidth]{figures/v5_eyz_profils.pdf}
\caption{Profil vertical de $|\epsilon_{yz}|$ à l'aplomb de son maximum : nul en surface, maximal
entre 80 et \SI{140}{mm}.}
\end{figure}
\begin{figure}[h]
\centering\includegraphics[width=\linewidth]{figures/v5_e1_cartes.pdf}
\caption{Déformation principale majeure en surface (demi-bogie $y>0$), équivalent de la figure 3-12
du TFE ; même échelle pour les six cartes.}
\end{figure}


\subsection{Tests unitaires}
\texttt{python -m pytest tests} (23 tests, dont 3 pour le contact et 5 pour le niveau 2 du module texture) : équivalence des solveurs dense et par rigidités
(3 fonds $\times$ 4 jeux d'interfaces, charges normales et tangentielles), symétrie $y$, équilibre
vertical global, et cohérence FFT/Hankel, déflexion absolue comprise.

%=====================================================================================================
\section{Mode d'emploi}
%=====================================================================================================
\subsection{Installation}
Python $\ge$ 3.9, NumPy, SciPy et Matplotlib ; pytest pour les tests. Aucune compilation.
\begin{lstlisting}
pip install numpy scipy matplotlib pytest
cd chausspec
python -m pytest tests          # vérification (environ 1 min)
python exemples/tfe_a340.py     # premier calcul
\end{lstlisting}

\subsection{Organisation du code}
\begin{center}\small
\begin{tabular}{p{3.6cm}p{11.2cm}}
\toprule
\texttt{materials.py} & lois : \texttt{Elastic}, \texttt{TwoS2P1D} (+ WLF, $\nu^*$), \texttt{GeneralizedKelvinVoigt}, \texttt{GeneralizedMaxwell}, \texttt{FrozenModulus}\\
\texttt{structure.py} & \texttt{Layer}, \texttt{Structure} (fond, interfaces)\\
\texttt{loads.py} & empreintes et spectres ; \texttt{Wheel}, \texttt{Loading}\\
\texttt{regimes.py} & \texttt{Static}, \texttt{Harmonic}, \texttt{Moving}\\
\texttt{kernel.py} & noyau : \texttt{StiffnessKernel} (défaut), \texttt{DenseKernel} (référence)\\
\texttt{spectral.py} & multiplicateurs spectraux de toutes les composantes\\
\texttt{grid.py} & \texttt{solve\_grid} (champs sur plans) et \texttt{GridResult}\\
\texttt{axisym.py} & \texttt{solve\_axisym} (Hankel)\\
\texttt{io.py}, \texttt{\_\_main\_\_.py} & cas décrits en JSON, écriture CSV/JSON/PNG\\
\addlinespace
\texttt{contact.py} & surfaces rainurées et aléatoires, MPD (ISO 13473-1), contact bande de roulement / texture (v0.3)\\
\texttt{texture.py} & couplage deux échelles, méthode héréditaire, réponse locale du massif homogène aux efforts normaux et tangentiels (v0.3--0.4)\\
\texttt{fe2d.py} & éléments finis 2D « pixel » en déformation plane généralisée : entaille des rainures, microstructure granulats / mortier (v0.4)\\
\bottomrule
\end{tabular}
\end{center}

\subsection{Utilisation en Python}
\begin{lstlisting}[language=Python]
import numpy as np, chausspec as cs

bbgb = cs.TwoS2P1D(E00=65, E0=30000, k=0.25, h=0.787, delta=1.58,
                   tau_ref=1.22, T_ref=9.3, nu=0.35)          # MPa, s
st = cs.Structure([cs.Layer(bbgb, 0.32, "BB+GB"),
                   cs.Layer(cs.Elastic(150, 0.35), 0.60, "GRH"),
                   cs.Layer(cs.Elastic(75, 0.35), 1.00, "Sol 1"),
                   cs.Layer(cs.Elastic(150, 0.35), 1.00, "Sol 2"),
                   cs.Layer(cs.Elastic(30000, 0.35), 2.00, "Substratum")],
                  bottom="rigid_smooth")

# empreinte hétérogène de l'Annexe III, normalisée à 370 kN
T = cs.GaussianPairs1D(P=[1.32, 0.80, 0.822], centers=[0.14846, 0.08804, 0.02864],
                       sig=[0.02824, 0.01695, 0.01741])      # transversal (y)
fp = cs.Separable(fx=cs.HalfEllipse1D(0.277), fy=T).scaled_to(370e3)
bogie = cs.Loading([cs.Wheel(fp, x0, y0) for x0 in (0, -1.98) for y0 in (-0.7, 0.7)])

res = cs.solve_grid(st, bogie, cs.Moving(0.66), depths=[0.0, 0.05, 0.32],
                    comps=("exx", "eyy", "ezz", "exy", "exz", "eyz"),
                    L=(16, 8), N=(1024, 1024), sym_y=True, window=(-3, 1.2, 0, 1.2))

print(res.argmax("eyy", 0.32))                    # (valeur, x, y) en base de couche liée
e1 = res.principal_strains(0.0)[..., 2]           # déformation principale majeure en surface
t, eps = res.time_signal("eyy", 0.32, x_gauge=0.0, y_gauge=0.7, speed=0.66)   # jauge
\end{lstlisting}

\subsection{Utilisation par fichier JSON}
\begin{lstlisting}
python -m chausspec exemples/cas_tfe_rectangulaire.json
\end{lstlisting}
Le fichier décrit la structure, les roues (\texttt{rect}, \texttt{circle}, \texttt{separable},
\texttt{map}), le régime (\texttt{static}, \texttt{moving}, \texttt{harmonic}), la grille et les
sorties : profondeurs, composantes et jauges. Le programme écrit dans le dossier de sortie un CSV
par champ (1\textsuperscript{re} ligne : $x$ ; 1\textsuperscript{re} colonne : $y$), un
\texttt{synthese.json} (extrema et positions, $\epsilon_1$ si les six déformations sont demandées),
les signaux de jauges et des cartes PNG.

\paragraph{Carte de pression mesurée} (\texttt{"type": "map"}) : fichier CSV séparé par
\texttt{;}, dont la 1\textsuperscript{re} ligne contient les abscisses $x$ des centres de pixels
(la 1\textsuperscript{re} case est ignorée) et la 1\textsuperscript{re} colonne les ordonnées $y$.
Le reste contient les pressions ; \texttt{unit} convertit en Pa (\texttt{1e6} pour des MPa) et
\texttt{force} renormalise la résultante. Voir \texttt{exemples/cas\_carte\_mesuree.json}.

\subsection{Conseils pratiques}
\begin{itemize}[nosep]
\item \textbf{Grandeurs en profondeur} (base de couche liée) : un pas de 3 à \SI{6}{cm} suffit ;
      calcul en quelques secondes.
\item \textbf{Grandeurs de surface} sous empreinte hétérogène : un pas $\le\SI{8}{mm}$ en
      transversal ; vérifier la convergence en doublant $N$.
\item \textbf{Mémoire} : chaque couple (composante, profondeur) stocke un spectre
      $N_y\times(N_x/2+1)$ ; \texttt{store="complex64"} divise la mémoire par deux et
      \texttt{window} réduit les sorties.
\item \textbf{Signe} : compression négative (contraintes et déformations). La charge roule vers
      $+x$ et la traînée viscoélastique est en $X<0$.
\item \textbf{Températures} : \texttt{TwoS2P1D(..., C1=, C2=, T=)} ; pour un gradient,
      découper l'enrobé en sous-couches.
\end{itemize}

%=====================================================================================================
\section{Limites et perspectives}\label{sec:limites}
%=====================================================================================================
\begin{itemize}[nosep]
\item \textbf{Inertie} : négligée. Pour le décollage ou l'atterrissage, il faudra ajouter les termes
      $-\rho\omega^2$ dans l'équation de Navier ($\omega=-k_1V$). La solution générale fait alors
      intervenir deux nombres d'onde verticaux (ondes P et S) ; la structure du code est inchangée.
\item \textbf{Interfaces partiellement collées} (ressorts de Goodman) : il suffit de modifier le
      raccordement, par une loi $T=k_s\,[U]$ à l'interface.
\item \textbf{Matériaux non linéaires} (GNT, sols) : hors de portée d'une méthode spectrale ; on
      peut les approcher par des modules sécants itérés couche par couche.
\item \textbf{Anisotropie} (transversalement isotrope) : faisable, mais la solution générale change.
\item \textbf{Perspectives thèse} : (1) brancher les cartes du prototype STAC (\texttt{PressureMap}) ;
      (2) étudier les critères de fissuration par le haut, avec $\epsilon_{yz}$ maximal en
      profondeur et les déformations principales près des nervures ; (3) mener des études de
      sensibilité massives sur la vitesse, la température et le type de pneu ; (4) rétrocalculer
      le HWD en fréquentiel avec le 2S2P1D (régime \texttt{Harmonic} et solveur axisymétrique) ;
      (5) ajouter des efforts tangentiels mesurés (freinage, virage), déjà pris en charge par le code.
\end{itemize}

%=====================================================================================================
\appendix
\section{Dérivation symbolique de la solution générale}\label{ann:sympy}
Le script \texttt{derivation/derive.py} injecte
$U=(a+b\,\xi z)e^{\mp\xi z}$ et $W=(c+d\,\xi z)e^{\mp\xi z}$ dans \eqref{eq:psv1}--\eqref{eq:psv2},
puis identifie les coefficients :
\begin{lstlisting}
famille e^{-xi z} : a = -c + kappa d ,  b = -d
famille e^{+xi z} : a =  c + kappa d ,  b =  d       (kappa = (lambda+3 mu)/(lambda+mu))
\end{lstlisting}
En réécrivant la seconde famille en fonction de $t=\xi(h-s)$, on obtient \eqref{eq:solgen}.

\section{Références}
\small
\begin{itemize}[nosep,leftmargin=*]
\item Burmister D.M. (1945). The general theory of stresses and displacements in layered systems. \emph{J. Applied Physics} 16, 89--94, 126--127, 296--302.
\item Chabot A., Chupin O., Deloffre L., Duhamel D. (2010). ViscoRoute 2.0: a tool for the simulation of moving load effects on asphalt pavement. \emph{Road Materials and Pavement Design} 11(2), 227--250.
\item Chupin O., Chabot A., Piau J.-M., Duhamel D. (2010). Influence of sliding interfaces on the response of a layered viscoelastic medium under a moving load. \emph{Int. J. Solids Struct.} 47, 3435--3446.
\item De Beer M., Fisher C. (2000, 2002). Stress-In-Motion (SIM) / VRSPTA : mesures tridimensionnelles des contraintes de contact pneu-chaussée. CSIR, Afrique du Sud.
\item Duhamel D., Chabot A., Tamagny P., Harfouche L. (2005). ViscoRoute : visco-elastic modeling for asphalt pavements. \emph{Bulletin des Laboratoires des Ponts et Chaussées} 258--259, 89--103.
\item Love A.E.H. (1929). The stress produced in a semi-infinite solid by pressure on part of the boundary. \emph{Phil. Trans. R. Soc. A} 228, 377--420.
\item Nakhaei M. (2020). PyMastic : multi-layered elastic analysis in Python. \url{https://github.com/Mostafa-Nakhaei/PyMastic}.
\item Olard F., Di Benedetto H. (2003). General ``2S2P1D'' model and relation between the linear viscoelastic behaviours of bituminous binders and mixes. \emph{Road Materials and Pavement Design} 4(2), 185--224.
\item Petitjean J. et al. (2001, 2002). Campagne PEP, aéroport de Toulouse-Blagnac (LCPC, Airbus).
\item Siddharthan R.V., Yao J., Sebaaly P.E. (1998). Pavement strain from moving dynamic 3D load distribution. \emph{J. Transportation Engineering} 124(6), 557--566.
\item Wang H., Al-Qadi I.L. (2010). Near-surface pavement failure under multiaxial stress state in thick asphalt pavement. \emph{Transportation Research Record} 2154, 91--99.
\end{itemize}
\normalsize
\noindent\emph{Les références bibliographiques sont à vérifier et à compléter dans Zotero avant toute diffusion.}

\end{document}
